\documentclass{ohpc-doc}

% Define Base OS and other local macros
\newcommand{\OS}{AlmaLinux}
\newcommand{\baseOS}{\OS{} 10}
\newcommand{\OSRepo}{\OS{}\_10}
\newcommand{\OSTree}{EL\_10}
\newcommand{\OSTag}{el10}
\newcommand{\baseos}{almalinux-10}
\newcommand{\baseosshort}{almalinux10}
\newcommand{\osimage}{almalinux:10}
\newcommand{\provisioner}{confluent}
\newcommand{\provheader}{\provisioner{}}
\newcommand{\rms}{SLURM}
\newcommand{\rmsshort}{slurm}
\newcommand{\arch}{x86\_64}

\newcommand{\isoimage}{AlmaLinux-10.0-x86\_64-dvd.iso}
\newcommand{\downloadpage}{https://almalinux.org/get-almalinux/}
\newcommand{\confluentprofile}{alma-10.0-x86\_64-default}

% Define package manager commands
\newcommand{\pkgmgr}{dnf}
\newcommand{\addrepo}{dnf -y config-manager --add-repo}
\newcommand{\chrootaddrepo}{nodeshell compute dnf -y config-manager --add-repo}
\newcommand{\clean}{dnf clean expire-cache}
\newcommand{\chrootclean}{dnf --installroot=\$CHROOT clean expire-cache}
\newcommand{\install}{dnf -y install}
\newcommand{\chrootinstall}{nodeshell compute dnf -y install}
\newcommand{\groupinstall}{dnf -y groupinstall}
\newcommand{\groupchrootinstall}{nodeshell compute dnf -y groupinstall}
\newcommand{\remove}{dnf -y remove}
\newcommand{\upgrade}{dnf -y upgrade}
\newcommand{\chrootupgrade}{dnf -y --installroot=\$CHROOT upgrade}
\newcommand{\tftppkg}{syslinux-tftpboot}
\newcommand{\beegfsrepo}{https://www.beegfs.io/release/beegfs\_7.2.1/dists/beegfs-rhel8.repo}
\newcommand{\cudarepo}{https://developer.download.nvidia.com/compute/cuda/repos/rhel10/x86\_64/cuda-rhel10.repo}

% boolean for os-specific formatting
\toggletrue{isCentOS}
\toggletrue{isCentOS_ww_slurm_x86}
\toggletrue{isSLURM}
\toggletrue{isx86}
\toggletrue{isCentOS_x86}

\begin{document}

% Introduction  --------------------------------------------------

\chapter{Introduction} \label{sec:introduction}
\input{common/install_header}
\input{common/intro}
\input{common/base_edition/edition}
\input{common/audience}
\input{common/requirements}
\input{common/inputs}

% Base Operating System --------------------------------------------

\chapter{Install Base Operating System (BOS)}
\input{common/bos}

% begin_ohpc_run
% ohpc_validation_newline
% ohpc_validation_comment Disable firewall
\begin{lstlisting}[language=bash,keywords={}]
[sms](*\#*) systemctl disable firewalld
[sms](*\#*) systemctl stop firewalld
\end{lstlisting}
% end_ohpc_run

% ------------------------------------------------------------------

\chapter{Install \Confluent{} and Provision Nodes with BOS} \label{sec:provision_compute_bos}
\input{common/confluent_compute_bos_intro}

\section{Enable \Confluent{} repository for local use} \label{sec:enable_confluent}
\input{common/enable_confluent_repo}

\section{Add provisioning services on {\em master} node} \label{sec:add_provisioning}
\input{common/install_provisioning_confluent_intro}

\section{Complete basic \Confluent{} setup for {\em master} node} \label{sec:setup_confluent}
\input{common/confluent_setup}

\section{Define {\em compute} image for provisioning}
\input{common/confluent_init_os_images_rhel}

\section{Add compute nodes into \Confluent{} database} \label{sec:confluent_add_nodes}
\input{common/add_confluent_hosts_intro}

\section{Boot compute nodes} \label{sec:boot_computes}
\input{common/reset_computes_confluent}

% begin_ohpc_run
% ohpc_validation_newline
% ohpc_validation_comment Check if provisioning is in progress
% ohpc_validation_newline
% ohpc_command while true; do
% ohpc_command    deployment_pending=false
% ohpc_command    while read line; do
% ohpc_command        dp=$(echo $line | cut -d ":" -f 2)
% ohpc_command          if [ "$dp" != " completed" ];
% ohpc_command          then
% ohpc_command               deployment_pending=true
% ohpc_command          fi
% ohpc_command    done <<< `nodedeploy compute`
% ohpc_command    if $deployment_pending;
% ohpc_command     then
% ohpc_command          echo "deployment still pending"
% ohpc_command          deployment_pending=false
% ohpc_command          sleep 60
% ohpc_command     else
% ohpc_command          break
% ohpc_command     fi
% ohpc_command done
% end_ohpc_run

\chapter{Install \OHPC{} Components} \label{sec:basic_install}
\input{common/install_ohpc_components_intro}

\section{Enable \OHPC{} repository for local use} \label{sec:enable_repo}
\input{common/enable_local_ohpc_repo_confluent}

% begin_ohpc_run
% ohpc_validation_newline
% ohpc_validation_comment Verify OpenHPC repository has been enabled before proceeding
% ohpc_validation_newline
% ohpc_command dnf repolist | grep -q OpenHPC
% ohpc_command if [ $? -ne 0 ];then
% ohpc_command    echo "Error: OpenHPC repository must be enabled locally"
% ohpc_command    exit 1
% ohpc_command fi
% end_ohpc_run

\input{common/almalinux_repos}

% begin_ohpc_run
\begin{lstlisting}[language=bash,keywords={},basicstyle=\fontencoding{T1}\fontsize{8.0}{10}\ttfamily,literate={ARCH}{\arch{}}1 {-}{-}1]
[sms](*\#*) (*\install*) epel-release

# Enable crb on sms
[sms](*\#*) /usr/bin/crb enable
\end{lstlisting}
% end_ohpc_run

Now \OHPC{} packages can be installed. To add the base package on the SMS
issue the following
% begin_ohpc_run
\begin{lstlisting}[language=bash,keywords={},basicstyle=\fontencoding{T1}\fontsize{8.0}{10}\ttfamily,literate={ARCH}{\arch{}}1 {-}{-}1]
[sms](*\#*)  (*\install*) ohpc-base
\end{lstlisting}
% end_ohpc_run

\input{common/automation}

\section{Setup time synchronization service on {\em master} node} \label{sec:add_ntp}
\input{common/time}

%\input{common/enable_pxe}

\section{Add resource management services on {\em master} node} \label{sec:add_rm}
\input{common/install_slurm}

\section{Optionally add \InfiniBand{} support services on {\em master} node} \label{sec:add_ofed}
\input{common/ibsupport_sms_centos}

\section{Optionally add \OmniPath{} support services on {\em master} node} \label{sec:add_opa}
\input{common/opasupport_sms_centos}

\subsection{Add \OHPC{} components} \label{sec:add_components}
\input{common/add_to_compute_confluent_intro}

% begin_ohpc_run
% ohpc_validation_comment Add OpenHPC components to compute instance
\begin{lstlisting}[language=bash,literate={-}{-}1,keywords={},upquote=true]
# Add Slurm client support meta-package
[sms](*\#*) (*\chrootinstall*) ohpc-slurm-client

# Add Network Time Protocol (NTP) support
[sms](*\#*) (*\chrootinstall*) chrony

# Add kernel drivers
[sms](*\#*) (*\chrootinstall*) kernel

# Enable crb
[sms](*\#*) nodeshell compute /usr/bin/crb enable

# Include nfs-utils
[sms](*\#*) (*\chrootinstall*) nfs-utils

# Include modules user environment
[sms](*\#*) (*\chrootinstall*) lmod-ohpc
\end{lstlisting}

% end_ohpc_run

% ohpc_comment_header Optionally add InfiniBand support services in compute node image \ref{sec:add_components}
% ohpc_command if [[ ${enable_ib} -eq 1 ]];then
% ohpc_indent 5
\begin{lstlisting}[language=bash,literate={-}{-}1,keywords={},upquote=true]
# Optionally add IB support and enable
[sms](*\#*) (*\groupchrootinstall*) "InfiniBand Support"
\end{lstlisting}
% ohpc_indent 0
% ohpc_command fi
% end_ohpc_run

\subsection{Customize system configuration} \label{sec:master_customization}
\input{common/confluent_customize}

% Additional commands when additional computes are requested

% begin_ohpc_run
% ohpc_validation_newline
% ohpc_validation_comment Update basic slurm configuration if additional computes defined
% ohpc_validation_comment This is performed on the SMS, nodes will pick it up config file is copied there later
% ohpc_command if [ ${num_computes} -gt 4 ];then
% ohpc_command    perl -pi -e "s/^NodeName=(\S+)/NodeName=${compute_prefix}[1-${num_computes}]/" /etc/slurm/slurm.conf
% ohpc_command    perl -pi -e "s/^PartitionName=normal Nodes=(\S+)/PartitionName=normal Nodes=${compute_prefix}[1-${num_computes}]/" /etc/slurm/slurm.conf
% ohpc_command fi
% end_ohpc_run

\subsection{Additional Customization ({\em optional})} \label{sec:addl_customizations}
\input{common/compute_customizations_intro}

\subsubsection{Increase locked memory limits}
\input{common/memlimits_confluent}

\subsubsection{Enable ssh control via resource manager}
\input{common/slurm_pam_confluent}

%\subsubsection{Add \Lustre{} client} \label{sec:lustre_client}
%\input{common/lustre-client}
%\input{common/lustre-client-centos-confluent}
%\input{common/lustre-client-post-confluent}

\subsubsection{Add GPU driver}
To add the optional NVIDIA GPU driver to the compute nodes, an
additional external \pkgmgr{} repository provided by NVIDIA
must be configured. Once the repository is configured the
GPU driver and the corresponding toolkit needs to be installed:

% begin_ohpc_run
% ohpc_validation_newline
% ohpc_command if [[ ${enable_nvidia_gpu_driver} -eq 1 ]];then
% ohpc_indent 5
\begin{lstlisting}[language=bash,literate={-}{-}1,keywords={},upquote=true]
# Add NVIDIA GPU driver repository
[sms](*\#*) nodeshell compute dnf -y config-manager --add-repo https://developer.download.nvidia.com/compute/cuda/repos/rhel9/x86_64/cuda-rhel9.repo
# Install the GPU driver
[sms](*\#*) nodeshell compute dnf -y module install nvidia-driver:latest-dkms
# Install the toolkit
[sms](*\#*) nodeshell compute dnf -y install cuda-toolkit
\end{lstlisting}
% ohpc_indent 0
% ohpc_command fi
% \end_ohpc_run


\subsubsection{Add \clustershell{}}
\input{common/clustershell}

\subsubsection{Add \genders{}}
\input{common/genders}

\subsubsection{Add Magpie}
\input{common/magpie}

\subsubsection{Add \conman{}} \label{sec:add_conman}
\input{common/conman}

\subsubsection{Add \nhc{}} \label{sec:add_nhc}
\input{common/nhc}
\input{common/nhc_slurm}

%\subsection{Identify files for synchronization} \label{sec:file_import}
%\input{common/import_confluent_files}
%\input{common/import_confluent_files_slurm}

%%%\subsection{Optional kernel arguments} \label{sec:optional_kargs}
%%%\input{common/conman_post}

\chapter{Install \OHPC{} Development Components}
\input{common/dev_intro}

\section{Development Tools} \label{sec:install_dev_tools}
\input{common/dev_tools}

\section{Compilers} \label{sec:install_compilers}
\input{common/compilers}

\section{MPI Stacks} \label{sec:mpi}
\input{common/mpi_slurm}

\section{Performance Tools} \label{sec:install_perf_tools}
\input{common/perf_tools}

\section{Setup default development environment}
\input{common/default_dev}

\section{3rd Party Libraries and Tools} \label{sec:3rdparty}
\input{common/third_party_libs_intro}

\input{common/third_party_libs}
\input{common/third_party_mpi_libs_x86}

\section{Optional Development Tool Builds} \label{sec:3rdparty_intel}
\input{common/oneapi_enabled_builds_slurm}

\chapter{Resource Manager Startup} \label{sec:rms_startup}
\input{common/slurm_startup_confluent}

\chapter{Run a Test Job} \label{sec:test_job}
\input{common/confluent_slurm_test_job}

\appendix

\input{common/automation_appendix}
\input{common/upgrade_confluent}
\input{common/test_suite}
\input{common/customization_appendix_centos}
\input{manifest}
\clearpage
\subsection{Package Signatures}
%\addcontentsline{toc}{section}{Appendix B - Package Signatures}

All of the RPMs provided via the \OHPC{} repository are signed with a GPG
signature. By default, the underlying package managers will verify these signatures during
installation to ensure that packages have not been altered. The RPMs can also
be manually verified and the public signing key fingerprint for the latest
repository is shown below: \\

\texttt{Fingerprint: 5E33 3CA3 A1BD BBC9 DF14 9D74 09AD FAE4 {\bf D722A692}} \\

\noindent The following command can be used to verify an RPM once it
has been downloaded locally by confirming if the package is signed, and if so,
indicating which key was used to sign it. The example below highlights usage
for a local copy of the \texttt{docs-ohpc} package and illustrates how the {\em
key ID} matches the fingerprint shown above.

\begin{lstlisting}[language=bash,keywords={}]
[sms](*\#*) rpm --checksig -v ohpc-release-4-1.el10.x86_64.rpm
ohpc-release-4-1.el10.x86_64.rpm:
    Header V3 RSA/SHA256 Signature, key ID d722a692: OK
    Header SHA256 digest: OK
    Header SHA1 digest: OK
    Payload SHA256 digest: OK
    V3 RSA/SHA256 Signature, key ID d722a692: OK
    MD5 digest: OK

\end{lstlisting}


\end{document}
